\section{Background}

\paragraph{Reproducibility in AI evaluation.}
AI evaluation commonly relies on public datasets, executable code, fixed metrics, and artifact checklists. Sensitive speech settings complicate this model because the evaluation evidence may include personally revealing acoustic or transcript-bearing content.

\paragraph{Sensitive speech data.}
Speech data can contain identity cues, operational details, medical or financial information, emotional content, and contextual signals that remain sensitive even after transcription. Releasing row-level examples can therefore create privacy and governance risks.

\paragraph{Claim-evidence alignment.}
For sensitive AI evaluation, the paper claim should specify its evidence layer, statistic, scope, limitation, privacy class, and validation status. This turns limitations into a claim-evidence alignment discipline rather than a defensive caveat. The framework is compatible with documentation-oriented responsible-AI practice such as dataset documentation \citep{datasheets} and risk-management framing \citep{nistairmf}, while making no legal-compliance claim.

\paragraph{Audit trails and operation records.}
Operation records preserve process-level reproducibility: what generator or validator produced an artifact, from which source inputs, under which version and environment, with which checksum. They do not expose local-only content, but they make the aggregate evidence trail inspectable.
